\documentclass[onecolumn,draftclsnofoot, 10pt, compsoc]{IEEEtran}

\usepackage{graphicx}
\usepackage[section]{placeins}
\usepackage{caption}

\usepackage{amssymb}                                         
\usepackage{amsmath}                                         
\usepackage{amsthm}                                

\usepackage{alltt}                                           
\usepackage{float}
\usepackage{color}
\usepackage{url}

\usepackage{balance}
\usepackage[TABBOTCAP, tight]{subfigure}
\usepackage{enumitem}
\usepackage{pstricks, pst-node}
\usepackage{url}
\usepackage{setspace}

\usepackage{etoolbox}
\AtBeginEnvironment{quote}{\singlespacing\vspace{-\topsep}\small}

%\input{pygments.tex}

\usepackage{geometry}
\geometry{left=0.75in,right=0.75in,top=0.75in,bottom=0.75in}
\parindent = 0.0 in
\parskip = 0.1 in

\usepackage{listings}
\lstset{language=C,caption={Descriptive Caption Text},label=DescriptiveLabel}
\lstdefinestyle{customc}{
	belowcaptionskip=1\baselineskip,
	breaklines=true,
	frame=L,
	xleftmargin=\parindent,
	language=C,
	showstringspaces=false,
	basicstyle=\footnotesize\ttfamily,
	keywordstyle=\bfseries\color{green!40!black},
	commentstyle=\itshape\color{purple!40!black},
	identifierstyle=\color{blue},
	stringstyle=\color{orange},
}

\lstdefinestyle{customasm}{
	belowcaptionskip=1\baselineskip,
	frame=L,
	xleftmargin=\parindent,
	language=[x86masm]Assembler,
	basicstyle=\footnotesize\ttfamily,
	commentstyle=\itshape\color{purple!40!black},
}


\def \ParSpace{\vspace{.75em}}
\def \GroupNumber{		17}
\def \Jeremy{			Jeremy Fischer}
\def \Class{		Operating Systems ii}
\def \School{	Oregon State University}
\def \Professor{		 Kevin McGrath}

\newcommand{\cred}[1]{{\color{red}#1}}
\newcommand{\cblue}[1]{{\color{blue}#1}}

\newcommand{\NameSigPair}[1]{
		\par
		\makebox[2.75in][r]{#1} \hfil 	\makebox[3.25in]{\makebox[2.25in]{\hrulefill} \hfill			
		\makebox[.75in]{\hrulefill}}
		\par\vspace{-12pt} \textit{
			\tiny\noindent
			\makebox[2.75in]{} \hfil		
			\makebox[3.25in]{
				\makebox[2.25in][r]{Signature} \hfill	\makebox[.75in][r]{Date}
			}
		}
}










%%%%%%%%%%%%%%%%%%%%%%%%%%%%%%%%%%%%%%%
\begin{document}
\begin{titlepage}
    \pagenumbering{gobble}
    \begin{singlespace}
    \includegraphics[height=4cm]{coe.eps}
        \hfill  
        \par\vspace{.2in}
        \centering
        \scshape{
            \vspace{.5in}
            \textbf{\Huge\Class}\par
            \large{
            	\today \\Fall Term
        	}
            \vfill
            {\large Prepared for}\par
            \huge \School\par
            \vspace{5pt}
            {\Large{\Professor}\par}
            {\large Prepared by }\par
           % Group\GroupNumber\par
            \vspace{5pt}
            {\Large
                {\Jeremy}\par
            }
            \vspace{20pt}
        }
        \begin{abstract}
        	This document explores file systems in the Linux, FreeBSD, and Windows operating systems. 
        	The document concludes with a comparison of the three's implementation.
        \end{abstract}     
    \end{singlespace}
\end{titlepage}
\newpage
\pagenumbering{arabic}
\tableofcontents
% 7. uncomment this (if applicable). Consider adding a page break.
%\listoffigures
%\listoftables
\clearpage









\section{Linux Implementation}
	In Linux, regular files are simply a stream of bytes.
	When an application requests to read a regular file, the stream of bytes are presented to the application along with an End Of File (EOF) condition.
	A regular file is referenced by the system by an inode.
	
	When a file is created, a unique number is associated to it.
	This is called the inode number.
	An inode holds information about the file such as \cite{llinuxInodeContents}:
	\begin{itemize}
		\item the file's permission mode
		\item the number of links in place for the file
		\item the file owners UID number
		\item the group GID number
		\item the file size represented in bytes
		\item the address of the datablocks
		\item the time the file was last modified
		\item the time any part of the inode was changed
	\end{itemize}
	
	When a file is not being used, its inode which resides on disk is called a disk inode.
	When the file is being used, the kernel places the inode onto a generic inode table and the inode now becomes a generic inode \cite{llinuxFiles}.
	
	In Linux, a directory is simply a file whose contents provide a mapping mechanism between the names of files in the directory and the file's actual datablocks.
	More specifically, a directory only holds the inode numbers and filenames of the files within it \cite{llinuxFiles}.
	
	\subsection{Ext4}
		For most Linux distributions the Ext file system is the default file system.
		Ext was created in 1992 and has since been revised three times leading to the Ext4, which was created in 2008.
		The Ext4 file system was designed to be backwards compatible. This means one can mount an Ext4 file system as Ext3.
		Ext4 includes newer features that reduce file fragmentation, allows for larger volumes and files, and uses delayed allocation to improve flash memory life.
		Ext4 is a journaling file system.
		This means before the operating makes a write to disk, metadata is written in a journal.
		This way if the power goes out in the middle of a disk write the operating system knows to finish the write from checking the journal.
		Without a journal the operating system isn't aware that it didn't finish a write, creating a possibility for corrupt data on disk. 
		When formatting an external drive that will be shared with other operating systems, one should not use Ext4 because Windows, macOS, and other devices can’t read Ext4 file systems \cite{llinuxEXT4}.
	
	\subsection{BtrFS}
		BtrFS, stands for “B-Tree File System.”
		BtrFS allows for drive pooling, on the fly snapshots, transparent compression, and online defragmentation \cite{llinuxEXT4}.
	
\section{FreeBSD Implementation}
	\subsection{UFS}
		UFS stands for "Unix File System."
		UFS is split into two layers: an upper layer (called the UFS) and a lower layer (called the FFS).
		The upper layer contains the file system's directories in an inode structure and metadata such as permissions and ownership.
		The lower layer provides data containers implemented as inodes \cite{freebsdUFS}.
		
		In UFS2, the upper and lower layer are extended to add 64-bit block pointers.
		This allows for volumes to grow up to 8 zettabytes, variable-sized blocks, extended flag fields, and additional 'birthtime' stamps. 
		Since FreeBSD 7.0, UFS also supports filesystem journaling \cite{freebsdUFS}.

	\subsection{ZFS}
		 ZFS supports a lot of advanced features including drive pooling, snapshots, and dynamic disk striping.
		 Each file in ZFS has a checksum, so the file system can tell if a file is corrupted or not \cite{llinuxEXT4}.
		 
		 ZFS is more than just a file system. 
		 ZFS combines the volume manager and the file system, which are traditionally kept separate.
		 The file system is now aware of the underlying structure of the disks. 
		 Traditional file systems could only be created on a single disk at a time. 
		 If a user wants to use two disks, then they have to mount two separate file systems. 
		 ZFS's combination of the volume manager and the file system solves this.
		 This enhancement also allows multiple file systems to share a pool of available storage. 
		 Due to ZFS's awareness of the physical layout, when an additional disk is added to the storage pool, each file system accessing that pool automatically grows \cite{freebsdZFS}.  

		
\section{Windows Implementation}
	\subsection{FAT}
		The FAT file system is characterized by the file allocation table (FAT).
		Two copies of the file allocation table are kept on the system with one acting as a backup in case one becomes corrupted. 
		A disk utilizing the FAT file system is organized in clusters.
		The size of each cluster is determined by the size of the overall volume. 
		When a new file is created it is given the first open location on the drive, and an entry is added in to the FAT table.
		This entry in the FAT table either indicates that this is the file's last cluster or it points to the file's next cluster.
		A big disadvantage of the FAT file system is updating the FAT table is time consuming, and doing so is essential \cite{windowsFAT}.
	
	\subsection{NTFS}
		NTFS organizes files into sorted directories.
		Unlike the FAT file system, NTFS is independent of the underlying hardware of the system, such as 512 byte sectors. 
		NTFS is also independent of any special locations on the disk, such as FAT tables.
		According to Microsoft \cite{windowsFAT}, the goals of NTFS are to provide:
		\begin{itemize}
			\item{
				reliability, which is especially desirable for high end systems and file servers}
			\item{
				a platform for added functionality}
			\item{
				support POSIX requirements}
			\item{
				removal of the limitations of the FAT and HPFS file systems}
		\end{itemize}

\section{Compare}
	Linux and FreeBSD are extremely similar in their underlying structure of directories and files.
	Because of this, most of the file systems supported by Linux are also supported by FreeBSD.
	Windows structures its file systems in a different manner.
	However, all three operating systems support the FAT file system.
	
	
\bibliographystyle{IEEEtran}
\bibliography{references.bib}

\end{document}